\documentclass[a4paper]{article}

\newcommand{\docTitle}{IT Architekturen}

\usepackage{datetime2}  % for dates
\usepackage[utf8]{inputenc} % UTF-8 input encoding
\usepackage[T1]{fontenc} % fixes \textquotedbl with pdflatex/OT1
\usepackage{textcomp}   % additional text symbols for typewriter/upquote
\usepackage{xcolor}     % for gray color
\usepackage{listings}
\usepackage{marginnote} % for horizontal margin notes
\usepackage{parskip}
\usepackage{amsmath,amssymb,mathtools}
\usepackage{everypage}
\usepackage{fancyhdr}
\usepackage{amsfonts}   % Mengen and natural numbers
\usepackage{enumerate}  % Custom Enumerations
\usepackage[inline]{enumitem}   % List spacing and labels
\usepackage{multicol}   % Multiple Columns
\usepackage{tabularx}

\usepackage{hyperref}   % Links
\hypersetup{
    pdftitle={\docTitle},
    pdfauthor={Moritz Köhler},
    pdfkeywords={DHBW, Hochschule},
    colorlinks=true,
    linkcolor=black,
    filecolor=magenta,
    urlcolor=cyan
}

% -------
% Images
% -------

\usepackage{import}
\usepackage{xifthen}
\usepackage{pdfpages}
\usepackage{tikz}
\usetikzlibrary{arrows.meta, positioning}

\newcommand{\incfig}[1]{%
    \def\svgwidth{\columnwidth}
    \import{./fig/}{#1.pdf_tex}
}

% ---------------------------------
% Vertical Side note (Bottom Left)
% ---------------------------------

\newcommand{\verticaldocinfo}{%
  \rotatebox{90}{\tiny\color{gray}\texttt{\DTMtoday\_\docTitle\_\jobname}}%
}

\AddEverypageHook{
  \put(-40pt,-650pt){\verticaldocinfo}%
}

% ---------------
% Math Shortcuts
% ---------------

% Number sets
\newcommand{\R}{\mathbb{R}}
\newcommand{\N}{\mathbb{N}}
\newcommand{\Z}{\mathbb{Z}}
\newcommand{\Q}{\mathbb{Q}}
\newcommand{\C}{\mathbb{C}}

% Vectors and matrices
\newcommand{\vect}[1]{\mathbf{#1}}
\newcommand{\mat}[1]{\mathbf{#1}}

% Derivatives
\newcommand{\dd}{\mathrm{d}}
\newcommand{\dpartial}[2]{\frac{\partial #1}{\partial #2}}
\newcommand{\ddt}{\frac{\dd}{\dd t}}

% Norms and absolute values
\newcommand{\norm}[1]{\left\lVert #1 \right\rVert}
\newcommand{\abs}[1]{\left\lvert #1 \right\rvert}

% Probability & Statistics
\newcommand{\E}{\mathbb{E}}
\newcommand{\Var}{\mathrm{Var}}
\newcommand{\Prob}{\mathbb{P}}

% Fields and Algebras
\newcommand{\K}{\mathbb{K}}
\newcommand{\F}{\mathbb{F}}

% Set Operations
\newcommand{\compl}[1]{#1^{\mathrm{c}}}

% Matrix and Vector Operations
\newcommand{\T}{^{\top}}
\newcommand{\inv}{^{-1}}
\newcommand{\conj}[1]{\overline{#1}}

% -------------------
% Main lecture macro 
% -------------------

\newcommand{\lecture}[1]{%
  \protect\marginnote{\protect\tiny\protect\color{gray}#1}
}

% -----------------------------
% Command for box with heading
% -----------------------------

\fboxsep2mm
\newcommand{\know}[2]{
    \noindent\fbox{
        \parbox{\textwidth-21pt}{

            \textbf{#1:}
            \vspace{0.125cm}

            #2
        }
    }
}

% -------------------
% Listings Languages
% -------------------

% Shared defaults for code listings
\lstset{
    basicstyle=\ttfamily\small,
    keywordstyle=\color{blue}\bfseries,
    commentstyle=\color{gray},
    stringstyle=\color{teal},
    numberstyle=\tiny\color{gray},
    numbers=left,
    stepnumber=1,
    numbersep=8pt,
    showstringspaces=false,
    frame=single,
    breaklines=true,
    tabsize=4,
    keepspaces=false,
    columns=flexible,
    upquote=true,
    extendedchars=false % disable extended char handling to avoid UTF-8 conflicts
}

% SQL syntax highlighting
\lstdefinelanguage{SQL}{
    morekeywords={SELECT,FROM,WHERE,INSERT,INTO,VALUES,UPDATE,SET,DELETE,CREATE,TABLE,PRIMARY,KEY,NOT,NULL,AND,OR,JOIN,LEFT,RIGHT,INNER,ON,AS,ORDER,BY,GROUP,HAVING,LIMIT},
    sensitive=false,
    morecomment=[l]{--},
    morestring=[b]'
}

% JSON syntax highlighting
\lstdefinelanguage{json}{
    morestring=[b]",
    stringstyle=\color{teal},
    comment=[l]{//},
    morecomment=[s]{/*}{*/},
        morekeywords={true,false,null},
    literate=
     *{0}{{{\color{blue}0}}}{1}
      {1}{{{\color{blue}1}}}{1}
      {2}{{{\color{blue}2}}}{1}
      {3}{{{\color{blue}3}}}{1}
      {4}{{{\color{blue}4}}}{1}
      {5}{{{\color{blue}5}}}{1}
      {6}{{{\color{blue}6}}}{1}
      {7}{{{\color{blue}7}}}{1}
      {8}{{{\color{blue}8}}}{1}
      {9}{{{\color{blue}9}}}{1}
            {:}{{{\color{black}:}}}{1}
            {,}{{{\color{black},}}}{1},
    showstringspaces=false
}

% Language specific styles
\lstdefinestyle{javastyle}{
    language=Java,
    basicstyle=\ttfamily\small,
    keywordstyle=\color{blue}\bfseries,
    commentstyle=\color{gray},
    stringstyle=\color{teal},
    numberstyle=\tiny\color{gray},
    numbers=left,
    stepnumber=1,
    numbersep=8pt,
    showstringspaces=false,
    frame=single,
    breaklines=true,
    tabsize=4,
    keepspaces=true,
    columns=fullflexible,
    upquote=true
}

\lstdefinestyle{sqlstyle}{
    language=SQL,
    basicstyle=\ttfamily\small,
    keywordstyle=\color{blue}\bfseries,
    commentstyle=\color{gray},
    stringstyle=\color{teal},
    numberstyle=\tiny\color{gray},
    numbers=left,
    stepnumber=1,
    numbersep=8pt,
    showstringspaces=false,
    frame=single,
    breaklines=true,
    tabsize=4,
    keepspaces=true,
    columns=fullflexible,
    upquote=true
}

\lstdefinestyle{pythonstyle}{
    language=Python,
    basicstyle=\ttfamily\small,
    keywordstyle=\color{blue}\bfseries,
    commentstyle=\color{gray},
    stringstyle=\color{teal},
    numberstyle=\tiny\color{gray},
    numbers=left,
    stepnumber=1,
    numbersep=8pt,
    showstringspaces=false,
    frame=single,
    breaklines=true,
    tabsize=4,
    keepspaces=true,
    columns=fullflexible,
    upquote=true
}

\lstdefinestyle{cstyle}{
    language=C,
    basicstyle=\ttfamily\small,
    keywordstyle=\color{blue}\bfseries,
    commentstyle=\color{gray},
    stringstyle=\color{teal},
    numberstyle=\tiny\color{gray},
    numbers=left,
    stepnumber=1,
    numbersep=8pt,
    showstringspaces=false,
    frame=single,
    breaklines=true,
    tabsize=4,
    keepspaces=true,
    columns=fullflexible,
    upquote=true
}

\lstdefinestyle{bashstyle}{
    language=bash,
    basicstyle=\ttfamily\small,
    keywordstyle=\color{blue}\bfseries,
    commentstyle=\color{gray},
    stringstyle=\color{teal},
    numberstyle=\tiny\color{gray},
    numbers=left,
    stepnumber=1,
    numbersep=8pt,
    showstringspaces=false,
    frame=single,
    breaklines=true,
    tabsize=4,
    keepspaces=true,
    columns=fullflexible,
    upquote=true
}

\lstdefinestyle{jsonstyle}{
    language=json,
    basicstyle=\ttfamily\small,
    keywordstyle=\color{blue}\bfseries,
    commentstyle=\color{gray},
    stringstyle=\color{teal},
    numberstyle=\tiny\color{gray},
    numbers=left,
    stepnumber=1,
    numbersep=8pt,
    showstringspaces=false,
    frame=single,
    breaklines=true,
    tabsize=4,
    keepspaces=true,
    columns=fullflexible,
    upquote=true
}

% Default listing style
\lstset{language={}}

\title{\docTitle}
\author{Moritz}
\date{\today}

\begin{document}
\maketitle
\tableofcontents

\lecture{2026-10-06}

\section{Organisatorisches}

Es gibt zwei Dozenten. Einer die erste Hälfte der Vorlesung, der andere die zweite Hälfte. Kennen sich aus IBM Zeiten.

Es gibt zwei Termine pro Woche, es müssen aber nicht beide stattfinden. Dies wird die Woche vorher geklärt.

In der Vorlesung geht es darum auch mal hinter den Vorhang zu gucken. Also wie Funktioniert Virtualisierung (Also zwei OS auf einem PC)

Ausführliche Vorstellungsrunde mit Firma, Vorerfahrungen, Erwartungen in einem Gespräch.

\subsection{Inhalte}

\begin{enumerate}
    \item Einführung
    \item Server Virtualisierung (Proxmox, Docker, Linux LXC, Typ 1 Hypervisor, x86 Virtualisierung)
    \item Zentralisierter Storage (SAN, NAS)
    \item Clusterarchitekturen (Datenbank-Cluster, Scale Out Data Center, KI Architekturen)
    \item IT Betrieb (Application-, System-, IT Service-Management)
\end{enumerate}

Netzwerk ist aufgrund der Zeit nicht drin.

\subsection{Prüfungsleistung}

Es wird eine Case Study, diese wird am 5. Termin (voraussichtlich 27.10) hochgeladen.

Bearbeitung in 3er Gruppen mit Referat 20min kurz vor der Klausurwoche mit Note. Bewertungsschema von der DHBW.

\section{Einführung}

\subsection{Einführung in IT-Architektur}

Was muss man so berücksichtigen:

\begin{itemize}
    \item Business Strategie
    \item Geschäftsprozesse und geschäftliche Fähigkeiten: Reiseziele (Katalog mit Bildern, ...), Buchen, Bezahlen
    \item Geschäftliche Anforderungen, Randbedingungen: Concurrent User, 1,5 TB Katalogdaten, Server Entfernung -> evtl. Daten Kopieren.
    \item Leistungsumfang deer IT Lösung: Erweiterung des Webshops, Anzeigen der Daten, Daten lokal, Prüfen der Verfügbarkeit nicht kopieren (da diese ja auch weg sein können), gibt es API?
    \item IT Infrastruktur und Betrieb: Storage, Netzwerkanbindung, Ausfallsicherung (Selbst, und Extern)
\end{itemize}

Es gibt Architekten und Enterprise Architekten, welche größere Entscheidungen festlegen (Warenkörbe für Softwarelizenzen, IT Sicherheit, Datenklassifikation)

Die meisten Architekten arbeiten an einem IT-System. Dort entweder als Anwendungs-Architektur (Funktionale Anforderungen, Benutzerschnittstellen) oder Infrastruktur-Architektur (nicht-funktionale Anforderungen, Netzwerk, operationelle Modelle)

\know{Schwerpunkt der Vorlesung:}{\dots ist die Infrastruktur-Architektur, also die nicht-funktionalen Anforderungen.}

Als Architekt arbeite ich immer mit dem Auftraggeber oder Bauherr!

\subsubsection{Vier Domänen des Architekturmodells nach TOGAF}

TOGAF (The Open Group Architecture Framework)

\begin{enumerate}
    \item Business Architecture Domain\\
    \hspace{5pt}
    \begin{itemize*}[itemjoin=\quad]
        \item Geschäftsprozesse (SCM, CRM, etc.)
        \item Organisation
        \item Portfolio
        \item Standorte
        \item Business Anforderungen
    \end{itemize*}
    \item Data Architecture Domain\\
    \hspace{5pt}
    \begin{itemize*}[itemjoin=\quad]
        \item Logisches Datenmodell
        \item Physisches Datenmodell
        \item Datenflüsse
    \end{itemize*}
    \item Application Architecture Domain\\
    \hspace{5pt}
    \begin{itemize*}[itemjoin=\quad]
        \item Logische Informationssystem
        \item Physisches Informationssystem
        \item Anwendungskoppelung (API, Messaging, BPM, etc.)
        \item Funktionale Anforderungen
    \end{itemize*}
    \item Technology Architecture Domain (unser Schwerpunkt)\\
    \hspace{5pt}
    \begin{itemize*}[itemjoin=\quad]
        \item Server
        \item Storage
        \item Netzwerk
        \item Virtualisierung
        \item Cluster
        \item Standards
        \item Data Center
        \item Middleware
        \item Tools
        \item IT Betriebsprozesse
        \item Nicht-funktionale Anforderungen (NFR)
    \end{itemize*}
\end{enumerate}

\subsubsection{TOGAF-Blume: Architecture Development Method (ADM)}

Es gibt auch die sogenannte \href{https://de.wikipedia.org/wiki/TOGAF#/media/Datei:TOGAF_ADM.svg}{TOGAF-Blume}.

Hier sollte man tatsächlich nur die relevanten Teile Nutzen.

Es entstehen im Laufe dieser Blume unterschiedliche Dokumente, die im Laufe der Entwicklung an unterschiedlichen Stellen Finalisiert werden. (S. 14)

\subsubsection{Architecture}

Hier gibt es zwei:

\begin{itemize}
    \item Architecture Definition Document (ADD)
    \item Architecture Requirements Specification (ARS)
\end{itemize}

Diese beiden Entstehen zum Start und sollten bei dem Migration Planning Finalisiert sein.

\begin{multicols*}{2}
\raggedcolumns

Bei ADD geht es um:

\begin{itemize}
    \item Leistungsfähigkeit
    \item Kapazität
    \item Verfügbarkeit 
    \item Sicherheit
    \item System Management
    \item Wartbarkeit
    \item Portabilität
    \item Skalierbarkeit
\end{itemize}
\columnbreak

Bei ARS geht es um:

\begin{itemize}
    \item Standorte
    \item Netzwerktopologie
    \item Benutzerzugänge
    \item Auf Logischem Level
    \item Und Physischem Level
\end{itemize}

\end{multicols*}

\lecture{}

\subsection{Dynamische IT Infrastrukturen}

\subsection{Cloud Computing}
% Mit Übungsaufgabe



\end{document}